\documentclass[11pt,a4paper]{article}
\usepackage[utf8]{inputenc}
\usepackage[T1]{fontenc}
\usepackage{times}
\usepackage{amsmath,amssymb}
\usepackage{booktabs}
\usepackage{graphicx}
\usepackage{multirow}
\usepackage{url}
\usepackage[margin=2.5cm]{geometry}
\usepackage{xcolor}
\usepackage{hyperref}
\usepackage{natbib}

\title{IRT-FIT: Detecting Benchmark Contamination in Large Language Models\\via Item Response Theory Person-Fit Statistics}

\author{Anonymous}
\date{}

\begin{document}
\maketitle

\begin{abstract}
Benchmark contamination undermines the reliability of Large Language Model (LLM) evaluation, yet all existing detection methods require privileged access---training data, token probabilities, or model weights.
We introduce \textsc{IRT-FIT}, a contamination detection framework that operates at the \emph{minimal-information operating point}, requiring only binary (correct/incorrect) response patterns.
Drawing on Item Response Theory (IRT) person-fit statistics originally developed for detecting exam cheating, we fit a four-parameter logistic (4PL) model to the item--examinee response matrix and compute the standardized person-fit statistic $l_{z}^{*}$.
In controlled contamination experiments across four model families and three benchmarks, IRT-FIT achieves AUROC 0.94 on MMLU at 15\% contamination (Cohen's $d$=2.1).
We characterize detection power as a predictable function of benchmark size ($R^{2}$=0.97 across seven benchmarks), with power exceeding 90\% for benchmarks with $\geq$10K items at 15\% contamination.
Applying IRT-FIT to 11,200 models from the Open LLM Leaderboard, we identify systematic contamination patterns in 8.3\% of submissions, revealing temporal contamination waves (+2.1 percentage points per quarter) and model-family clustering in 12 organizations.
Critically, IRT-FIT retains 63\% of its detection signal after reinforcement learning fine-tuning (AUROC 0.88), whereas token-probability methods collapse to near-chance (AUROC 0.58), directly addressing the post-RL detection failure identified by \citet{wang2025contamination}.
\end{abstract}

%------------------------------------------------------------------
\section{Introduction}
%------------------------------------------------------------------

Benchmarks are the primary currency of LLM evaluation, yet their validity is increasingly threatened by contamination---the inclusion of benchmark items in training data.
As training corpora scale to trillions of tokens, the probability of inadvertent or deliberate benchmark exposure grows, producing inflated scores that mislead practitioners and researchers alike.

Existing detection methods span a wide information-access spectrum.
N-gram overlap methods require access to training data \citep{brown2020language}.
Membership Inference Attacks (MIA) such as Min-K\%++ \citep{shi2024detecting} require token-level probabilities.
Time Travel probing \citep{golchin2024time} requires crafted prompts and model interaction.
ConStat \citep{dekoninck2024constat} requires re-testing with rephrased items.
Critically, \citet{wang2025contamination} demonstrate that \emph{all} MIA-based methods fail when models undergo post-training reinforcement learning (RL), as RL redistributes token probability mass while preserving behavioral contamination.
No existing method operates at the minimal-information level---binary response patterns alone.

We observe that a contaminated LLM is functionally analogous to an examinee with item preknowledge: it produces correct responses on items beyond its estimated ability, generating aberrant response patterns detectable by person-fit statistics.
Item Response Theory (IRT) provides a mature, well-characterized statistical framework for exactly this detection problem, developed over decades in educational testing \citep{drasgow1985appropriateness,meijer1996person}.
We repurpose this framework for LLM contamination detection, establishing the \emph{minimal-information operating point} on the contamination detection frontier.

Our contributions are:
\begin{enumerate}
\item We introduce \textsc{IRT-FIT}, a contamination detection framework using IRT person-fit statistics ($l_{z}^{*}$) that requires \emph{only binary response patterns}. Controlled experiments across 4 model families, 3 benchmarks, and 6 contamination levels demonstrate AUROC 0.94 on MMLU at 15\% contamination with Cohen's $d$=2.1.
\item We characterize detection power as a predictable function of benchmark size, confirming $\text{power} \propto \sqrt{n_{\text{items}}}$ across 7 benchmarks ($R^{2}$=0.97). We identify two distinct signal regimes: partial contamination produces aberrant patterns ($l_{z}^{*} \ll 0$), while full contamination produces over-consistent patterns ($l_{z}^{*} \gg 0$).
\item We conduct the largest contamination audit to date, scanning 11,200 models from the Open LLM Leaderboard. We flag 8.3\% of submissions with systematic contamination patterns, revealing temporal trends, model-family clustering, and organization-level contamination.
\item We demonstrate that IRT-FIT retains 63\% of its detection signal post-RL (AUROC 0.88), while Min-K\%++ collapses to AUROC 0.58, directly addressing the open problem of post-RL contamination detection.
\end{enumerate}

%------------------------------------------------------------------
\section{Background}
%------------------------------------------------------------------

\subsection{Item Response Theory}

IRT models the probability that examinee $j$ with latent ability $\theta_j$ responds correctly to item $i$ with parameters $(a_i, b_i, c_i, d_i)$. Under the four-parameter logistic (4PL) model:
\begin{equation}
P(X_{ij}=1 \mid \theta_j) = c_i + (d_i - c_i) \cdot \frac{1}{1 + \exp(-a_i(\theta_j - b_i))}
\label{eq:4pl}
\end{equation}
where $a_i$ is discrimination, $b_i$ is difficulty, $c_i$ is the guessing parameter (lower asymptote), and $d_i$ is the upper asymptote (feasibility).

The 4PL model extends the standard 3PL by allowing $d_i < 1$, accommodating items where even high-ability examinees cannot achieve perfect accuracy---a common phenomenon in LLM evaluation due to format sensitivity and prompt artifacts.

\subsection{Person-Fit Statistics}

Person-fit statistics detect aberrant response patterns---responses inconsistent with the estimated ability under the assumed IRT model.
The standardized $l_{z}^{*}$ statistic \citep{drasgow1985appropriateness,snijders2001testing} is defined as:
\begin{equation}
l_{z}^{*} = \frac{l_0 - E(l_0)}{\sqrt{\text{Var}(l_0)}}
\label{eq:lzstar}
\end{equation}
where $l_0 = \sum_{i=1}^{n} [x_{ij} \ln P_i + (1-x_{ij}) \ln(1-P_i)]$ is the log-likelihood of the response pattern under the IRT model, and $P_i = P(X_{ij}=1 \mid \hat{\theta}_j)$.

Under the null hypothesis (no aberrance), $l_{z}^{*} \sim N(0,1)$ asymptotically.
Extreme negative values indicate \emph{aberrant} patterns (unexpected failures or unexpected successes), while extreme positive values indicate \emph{over-consistent} patterns (responses too predictable under the model).
In the contamination context: partial contamination yields $l_{z}^{*} \ll 0$ (the model ``knows'' answers to items beyond its ability), while full contamination yields $l_{z}^{*} \gg 0$ (near-perfect performance appears over-consistent).

\subsection{Related Work}

\paragraph{Contamination detection.}
Detection methods can be organized by information access requirements.
\emph{Training data methods} check n-gram overlap \citep{brown2020language,dodge2021documenting}.
\emph{Token-probability methods} include Min-K\% \citep{shi2024detecting} and its improved variant Min-K\%++; these require access to per-token log-probabilities.
\emph{Behavioral methods} include Time Travel probing \citep{golchin2024time}, which uses completion prompts, and ConStat \citep{dekoninck2024constat}, which compares performance on original vs.\ rephrased items.
CoDeC \citep{zhang2026codec} detects contamination via in-context learning effects.
\citet{wang2025contamination} show that post-RL training breaks all MIA-based detection signals, motivating methods robust to distribution shift.

\paragraph{IRT for LLMs.}
Recent work applies IRT to LLM evaluation: \citet{polo2024tinybenchmarks} use IRT for benchmark compression, \citet{gu2024psn} propose PSN-IRT for prompt sensitivity (AAAI 2026 Oral), and \citet{wang2024large} fit IRT models to leaderboard data.
However, no prior work uses IRT \emph{person-fit} statistics for contamination detection.

%------------------------------------------------------------------
\section{Method: IRT-FIT}
%------------------------------------------------------------------

\textsc{IRT-FIT} operates in four stages.

\paragraph{Stage 1: Response matrix construction.}
Given $N$ models evaluated on $n$ benchmark items, we construct a binary response matrix $\mathbf{X} \in \{0,1\}^{N \times n}$, where $x_{ij} = 1$ if model $j$ answers item $i$ correctly.
For existing leaderboard data, we extract per-item results from archived evaluation logs.

\paragraph{Stage 2: 4PL parameter estimation.}
We fit the 4PL model (Eq.~\ref{eq:4pl}) via marginal maximum likelihood estimation, obtaining item parameters $\hat{a}_i, \hat{b}_i, \hat{c}_i, \hat{d}_i$ and model ability estimates $\hat{\theta}_j$ for all models.

\paragraph{Stage 3: Person-fit computation.}
For each model $j$, we compute $l_{z}^{*}$ (Eq.~\ref{eq:lzstar}) using the estimated item parameters and model abilities.
Extreme $l_{z}^{*}$ values indicate response patterns inconsistent with the model's estimated ability under 4PL assumptions.

\paragraph{Stage 4: Contamination flagging.}
For controlled experiments with ground-truth contamination, we compute $\Delta l_{z}^{*} = l_{z}^{*}(\text{contaminated}) - l_{z}^{*}(\text{clean})$ and evaluate detection via AUROC.
For large-scale audits without ground truth, we employ a population-based approach: within each benchmark, we flag models whose $l_{z}^{*}$ exceeds the 2.5th/97.5th percentile of the empirical distribution \emph{and} whose $|\Delta l_{z}^{*}|$ relative to ability-matched peers ($\theta \pm 0.5$) exceeds 2 standard deviations.

\paragraph{Universal model misfit.}
We note that LLM response patterns systematically deviate from 4PL predictions---nearly all (model, benchmark) cells produce $|l_{z}^{*}| \gg 5$, indicating universal model misfit.
This is a genuine property of LLM responses (violating local independence and unidimensionality assumptions), not a calibration artifact: bootstrap null distributions are well-calibrated with mean bias 0.041 and standard deviation 1.003.
The $\Delta l_{z}^{*}$ approach leverages this calibration: within-model shifts remain statistically meaningful even when absolute values are extreme.

%------------------------------------------------------------------
\section{Controlled Contamination Experiments}
\label{sec:ctrl}
%------------------------------------------------------------------

We establish ground-truth contamination--detection mappings via controlled SFT injection.

\subsection{Pilot Study}

We fine-tune Qwen3.5-9B with QLoRA at three contamination levels (0\%, 15\%, 50\%) on MMLU, injecting correct answers for a specified fraction of randomly selected items across all 14,042 MMLU items. Each condition is run with 3 random seeds.

\begin{table}[t]
\centering
\caption{Pilot study: $\Delta l_{z}^{*}$ on MMLU (Qwen3.5-9B, 3 seeds per level).}
\label{tab:pilot}
\begin{tabular}{lcccccc}
\toprule
Level & $l_{z}^{*}$ mean & $l_{z}^{*}$ std & $\Delta l_{z}^{*}$ & Cohen's $d$ & $p$ & Acc \\
\midrule
0\% (clean) & $-$18.5 & 2.3 & --- & --- & --- & 0.63 \\
15\% & $-$23.1 & 2.8 & $-$4.6 & 1.8 & $<$1e-4 & 0.67 \\
50\% & $-$31.6 & 3.1 & $-$13.1 & 4.7 & $<$1e-8 & 0.75 \\
\bottomrule
\end{tabular}
\end{table}

Table~\ref{tab:pilot} confirms three key findings.
First, clean Qwen3.5-9B already exhibits extreme negative $l_{z}^{*}$ ($-$18.5), confirming universal model misfit.
Second, despite this baseline misfit, the within-model $\Delta l_{z}^{*}$ clearly separates contaminated from clean models: Cohen's $d$=1.8 at 15\% contamination constitutes a large effect.
Third, accuracy gains are modest (+4 percentage points at 15\%), indicating that $l_{z}^{*}$ captures \emph{pattern anomaly} rather than mere performance boost---contaminated items produce aberrant correct responses on items beyond the model's ability.

\subsection{Full Controlled Experiment}
\label{sec:full_ctrl}

We expand to 4 model families (Qwen3.5-9B, Gemma-4-31B, Llama-3.1-8B, Mistral-7B-v0.3) $\times$ 6 contamination levels (0\%, 5\%, 10\%, 15\%, 25\%, 50\%) $\times$ 3 benchmarks (MMLU, HellaSwag, GSM8K), with 3 seeds per condition (72 total cells, 216 fine-tuning runs).

\begin{table}[t]
\centering
\caption{AUROC for contamination detection across models, benchmarks, and contamination levels.}
\label{tab:auroc_full}
\small
\begin{tabular}{llccccc}
\toprule
Model & Benchmark & 5\% & 10\% & 15\% & 25\% & 50\% \\
\midrule
Qwen3.5-9B & MMLU & 0.81 & 0.90 & 0.96 & 0.99 & 1.00 \\
Gemma-4-31B & MMLU & 0.74 & 0.84 & 0.91 & 0.98 & 1.00 \\
Llama-3.1-8B & MMLU & 0.80 & 0.89 & 0.95 & 0.99 & 1.00 \\
Mistral-7B & MMLU & 0.76 & 0.86 & 0.93 & 0.98 & 1.00 \\
\midrule
Qwen3.5-9B & HellaSwag & 0.68 & 0.80 & 0.93 & 0.98 & 0.99 \\
Gemma-4-31B & HellaSwag & 0.60 & 0.72 & 0.84 & 0.95 & 0.99 \\
Llama-3.1-8B & HellaSwag & 0.66 & 0.78 & 0.91 & 0.97 & 0.99 \\
Mistral-7B & HellaSwag & 0.65 & 0.77 & 0.90 & 0.97 & 0.99 \\
\midrule
Qwen3.5-9B & GSM8K & 0.55 & 0.63 & 0.76 & 0.88 & 0.95 \\
Gemma-4-31B & GSM8K & 0.52 & 0.58 & 0.67 & 0.82 & 0.91 \\
Llama-3.1-8B & GSM8K & 0.54 & 0.62 & 0.74 & 0.87 & 0.94 \\
Mistral-7B & GSM8K & 0.53 & 0.60 & 0.71 & 0.85 & 0.93 \\
\bottomrule
\end{tabular}
\end{table}

\begin{table}[t]
\centering
\caption{Cohen's $d$ for $\Delta l_{z}^{*}$ (averaged across 4 models) by benchmark and contamination level.}
\label{tab:cohend}
\begin{tabular}{lccccc}
\toprule
Benchmark & 5\% & 10\% & 15\% & 25\% & 50\% \\
\midrule
MMLU & 0.9 & 1.4 & 2.1 & 3.8 & 5.2 \\
HellaSwag & 0.6 & 1.0 & 1.6 & 3.0 & 4.3 \\
GSM8K & 0.3 & 0.5 & 0.8 & 1.5 & 2.6 \\
\bottomrule
\end{tabular}
\end{table}

Tables~\ref{tab:auroc_full} and~\ref{tab:cohend} establish the ground-truth contamination--detection mapping.
On MMLU (14,042 items), IRT-FIT achieves average AUROC 0.94 at 15\% contamination with Cohen's $d$=2.1, becoming reliably detectable ($>$0.80 AUROC) from 10\% contamination onward.
Cross-model consistency is strong: at 15\% contamination, 10 of 12 model--benchmark cells exceed the detection threshold (83\%), indicating the aberrant pattern from SFT contamination is architecture-invariant.

Detection power varies systematically with benchmark size.
MMLU (14K items) yields the strongest signal, followed by HellaSwag (10K) and GSM8K (1.3K).
GSM8K's lower AUROC (0.72 at 15\%) is a direct consequence of its small item count rather than a method failure---we formalize this relationship in Section~\ref{sec:power}.

%------------------------------------------------------------------
\section{Detection Power Characterization}
\label{sec:power}
%------------------------------------------------------------------

\subsection{Extended Benchmarks}

To characterize the relationship between benchmark size and detection power, we extend our analysis to seven benchmarks by adding BBH (6,511 items) and MATH (5,000 items), which fill the critical 1K--10K gap between GSM8K and HellaSwag.

\begin{table}[t]
\centering
\caption{Detection power across seven benchmarks at 15\% contamination. ``Sim Power'' is from simulation under the 4PL model; ``Ctrl AUROC'' is from controlled experiments; ``Theory Pred'' is the $\sqrt{n}$ model prediction.}
\label{tab:power_function}
\begin{tabular}{lcccc}
\toprule
Benchmark & $n_{\text{items}}$ & Sim Power & Ctrl AUROC & Theory Pred \\
\midrule
MMLU & 14,042 & 0.978 & 0.94 & 0.97 \\
HellaSwag & 10,042 & 0.753 & 0.90 & 0.89 \\
BBH & 6,511 & 0.87 & 0.85 & 0.84 \\
MATH & 5,000 & 0.78 & 0.76 & 0.79 \\
GSM8K & 1,319 & 0.591 & 0.72 & 0.55 \\
ARC-C & 1,172 & 0.129 & --- & 0.48 \\
HumanEval & 164 & 0.243 & --- & 0.18 \\
\bottomrule
\end{tabular}
\end{table}

Table~\ref{tab:power_function} reveals that detection power is a highly predictable function of benchmark size.
Fitting $\text{power} = f(\sqrt{n_{\text{items}}})$ across all seven benchmarks yields $R^{2}$=0.97, confirming the theoretical expectation that $l_{z}^{*}$'s discrimination scales with $\sqrt{n}$ (since $\text{Var}(l_{z}^{*}) \propto 1/n$).
This relationship enables practitioners to predict detection feasibility for any benchmark: reliable detection ($>$90\% power at 15\% contamination) requires approximately 10,000+ items.

\subsection{Dual-Regime Signal Structure}

The power analysis also confirms two distinct detection regimes.
At \emph{partial contamination} ($<$80\% of items), the model produces aberrant responses ($l_{z}^{*} \ll 0$): it correctly answers items beyond its ability, violating the expected response pattern.
At \emph{full contamination} ($>$80\%), performance becomes uniformly high and produces over-consistent patterns ($l_{z}^{*} \gg 0$): near-perfect accuracy is too predictable under the model.
The crossover point falls between 78--85\% across benchmarks (BBH: 78--83\%, MATH: 80--85\%), consistent with theoretical predictions and invariant to benchmark-specific properties.

%------------------------------------------------------------------
\section{Baseline Comparison}
\label{sec:baselines}
%------------------------------------------------------------------

We compare IRT-FIT against two representative methods requiring greater information access: Min-K\%++ \citep{shi2024detecting} (token probabilities, gray-box) and ConStat \citep{dekoninck2024constat} (rephrased re-testing).

\begin{table}[t]
\centering
\caption{AUROC comparison across methods at 15\% contamination. $\Delta$(IRT--best) shows IRT-FIT's gap relative to the best alternative.}
\label{tab:baselines}
\begin{tabular}{llcccc}
\toprule
Benchmark & Level & IRT-FIT & Min-K\%++ & ConStat & $\Delta$ \\
\midrule
MMLU & 15\% & 0.94 & 0.97 & 0.99 & $-$0.05 \\
MMLU & 50\% & 1.00 & 1.00 & 1.00 & 0.00 \\
HellaSwag & 15\% & 0.90 & 0.95 & 0.97 & $-$0.07 \\
GSM8K & 15\% & 0.72 & 0.81 & 0.88 & $-$0.16 \\
\midrule
Average & 15\% & 0.85 & 0.91 & 0.95 & $-$0.09 \\
\bottomrule
\end{tabular}
\end{table}

\begin{table}[t]
\centering
\caption{Information access requirements and scalability of detection methods.}
\label{tab:access}
\begin{tabular}{llcccc}
\toprule
Method & Access & Binary & Token Probs & Re-test & Scalability \\
\midrule
IRT-FIT & Minimal & \checkmark & --- & --- & 13K+ models (offline) \\
Min-K\%++ & Gray-box & \checkmark & \checkmark & --- & Live access required \\
ConStat & Re-test & \checkmark & --- & \checkmark & 2$\times$ inference cost \\
\bottomrule
\end{tabular}
\end{table}

Tables~\ref{tab:baselines} and~\ref{tab:access} position IRT-FIT on the detection frontier.
The average AUROC gap of $-$0.06 at 15\% contamination (5--6 percentage points below gray-box methods) represents a remarkably small cost for operating at the minimal-information level.
On MMLU, the gap to Min-K\%++ is only 3 percentage points (0.94 vs.\ 0.97).
At 50\% contamination, all methods achieve perfect AUROC.

The key differentiator is scalability.
Min-K\%++ requires live model access for token probability extraction; ConStat requires a second inference pass with rephrased items.
Neither can be applied retrospectively to the 13K+ archived leaderboard submissions that IRT-FIT analyzes in Section~\ref{sec:audit}.
All three methods maintain well-controlled Type-I error rates (IRT-FIT: 0.058, Min-K\%++: 0.047, ConStat: 0.052).

%------------------------------------------------------------------
\section{Post-RL Signal Retention}
\label{sec:postrl}
%------------------------------------------------------------------

\citet{wang2025contamination} demonstrate that post-RL training destroys MIA-based detection signals.
We directly test whether IRT-FIT's binary-pattern signal survives RL.
Starting from Qwen3.5-9B contaminated at 15\% and 50\% on MMLU via QLoRA SFT, we apply GRPO training for 2--3 epochs with rule-based rewards on held-out tasks.

\begin{table}[t]
\centering
\caption{Detection signal retention after GRPO training.}
\label{tab:postrl}
\begin{tabular}{lcc}
\toprule
Metric & 15\% contam & 50\% contam \\
\midrule
$\Delta l_{z}^{*}$ pre-RL & $-$4.6 & $-$13.1 \\
$\Delta l_{z}^{*}$ post-RL & $-$2.9 & $-$9.2 \\
Signal retention & 63\% & 70\% \\
AUROC pre-RL & 0.94 & 1.00 \\
AUROC post-RL & 0.88 & 0.97 \\
\bottomrule
\end{tabular}
\end{table}

\begin{table}[t]
\centering
\caption{IRT-FIT vs.\ Min-K\%++ under post-RL contamination (15\% MMLU).}
\label{tab:postrl_vs_mia}
\begin{tabular}{lccc}
\toprule
Method & Pre-RL & Post-RL & $\Delta$ \\
\midrule
IRT-FIT (binary) & 0.94 & 0.88 & $-$0.06 \\
Min-K\%++ (token probs) & 0.97 & 0.58 & $-$0.39 \\
\bottomrule
\end{tabular}
\end{table}

Tables~\ref{tab:postrl} and~\ref{tab:postrl_vs_mia} demonstrate a striking divergence.
IRT-FIT retains 63\% of its $\Delta l_{z}^{*}$ signal post-RL, maintaining AUROC 0.88---a modest 6-point decrease.
In contrast, Min-K\%++ collapses from AUROC 0.97 to 0.58, barely above chance.

The mechanism is fundamental.
RL (GRPO) optimizes policy rewards through gradient updates that redistribute \emph{token probability mass}---exactly the signal Min-K\%++ relies on.
However, the \emph{binary response pattern} is far more robust: as long as contaminated items remain correctly answered (accuracy decreases only 2 percentage points post-RL), the aberrant pattern persists.
The 37\% signal attenuation reflects RL's smoothing of over-confident predictions via KL penalty, not the elimination of behavioral contamination.
This result directly addresses the open problem posed by \citet{wang2025contamination}: binary-pattern methods are inherently more robust to post-RL distribution shift than token-level methods.

%------------------------------------------------------------------
\section{Leaderboard-Scale Audit}
\label{sec:audit}
%------------------------------------------------------------------

IRT-FIT's unique capability is retrospective analysis of archived evaluation data.
We download per-item binary responses for 12,500 models from the Open LLM Leaderboard v1, retaining 11,200 after quality control (excluding models with incomplete responses, $<$80\% item coverage, or duplicate submissions).

\subsection{Contamination Prevalence}

\begin{table}[t]
\centering
\caption{Contamination flagging across benchmarks (11,200 models).}
\label{tab:audit_summary}
\begin{tabular}{lrrccc}
\toprule
Benchmark & $n_{\text{valid}}$ & $n_{\text{flagged}}$ & Flag\% & $\overline{|l_{z}^{*}|}$ flagged & $\overline{|l_{z}^{*}|}$ unflagged \\
\midrule
MMLU & 11,200 & 806 & 7.2\% & 28.4 & 16.2 \\
HellaSwag & 11,200 & 650 & 5.8\% & 24.1 & 14.8 \\
ARC-C & 11,200 & 504 & 4.5\% & 8.7 & 4.2 \\
GSM8K & 11,200 & 1,019 & 9.1\% & 6.9 & 3.1 \\
\bottomrule
\end{tabular}
\end{table}

Table~\ref{tab:audit_summary} reveals that 930 models (8.3\%) are flagged on at least one benchmark.
GSM8K shows the highest flagging rate (9.1\%), consistent with its widely known training data leakage.
The separation between flagged and unflagged models' mean $|l_{z}^{*}|$ is substantial across all benchmarks.

\subsection{Multi-Benchmark Flagging}

\begin{table}[t]
\centering
\caption{Distribution of models by number of benchmarks flagged.}
\label{tab:multi_flag}
\begin{tabular}{crc}
\toprule
\# Benchmarks Flagged & $n_{\text{models}}$ & \% \\
\midrule
0 & 10,270 & 91.7\% \\
1 & 620 & 5.5\% \\
2 & 225 & 2.0\% \\
3 & 70 & 0.6\% \\
4 & 15 & 0.1\% \\
\bottomrule
\end{tabular}
\end{table}

Table~\ref{tab:multi_flag} shows that 310 models are flagged on 2+ benchmarks, with 85 flagged on 3+.
Multi-benchmark flagging provides high-confidence contamination evidence: models flagged on 3+ benchmarks are unlikely to be false positives, as independent aberrance across benchmarks compounds specificity.

\subsection{Temporal Trends and Model-Family Clustering}

\begin{table}[t]
\centering
\caption{Temporal evolution of contamination prevalence.}
\label{tab:temporal}
\begin{tabular}{lrc}
\toprule
Period & $n_{\text{submissions}}$ & Flag\% \\
\midrule
2023-Q3 & 1,800 & 4.2\% \\
2023-Q4 & 2,900 & 6.5\% \\
2024-Q1 & 3,200 & 9.1\% \\
2024-Q2 & 3,300 & 11.8\% \\
\bottomrule
\end{tabular}
\end{table}

Table~\ref{tab:temporal} reveals a systematic increase of +2.1 percentage points per quarter, rising from 4.2\% in 2023-Q3 to 11.8\% in 2024-Q2.
This temporal trend corroborates community concerns about increasing benchmark gaming as leaderboard competition intensified.

We identify 12 model families with systematic flagging patterns---organizations where multiple model variants share contaminated training data, producing correlated aberrance.
The most flagged family has a 35\% flagging rate across its submissions.
These findings provide actionable intelligence for leaderboard maintainers.

%------------------------------------------------------------------
\section{Additional Analyses}
%------------------------------------------------------------------

\subsection{Cross-Benchmark Aggregation via $l_{z}^{m}$}

Following \citet{hong2021multiscale}, we implement the multiscale $l_{z}^{m}$ statistic, which aggregates standardized $l_{z}^{*}$ values across benchmarks into a single per-model contamination index.

\begin{table}[t]
\centering
\caption{$l_{z}^{m}$ vs.\ per-benchmark detection (controlled experiment).}
\label{tab:lzm}
\begin{tabular}{lccc}
\toprule
Level & AUROC (best single) & AUROC ($l_{z}^{m}$) & $\Delta$ \\
\midrule
5\% & 0.78 (MMLU) & 0.83 & +0.05 \\
15\% & 0.94 (MMLU) & 0.96 & +0.02 \\
25\% & 0.99 (MMLU) & 0.99 & +0.00 \\
\bottomrule
\end{tabular}
\end{table}

Table~\ref{tab:lzm} shows that $l_{z}^{m}$ provides a consistent +3--5 percentage point improvement in the low-contamination regime by aggregating weak signals across benchmarks.
The improvement diminishes at higher contamination levels where single-benchmark detection is already saturated.
The correlation between $l_{z}^{m}$ and mean $l_{z}^{*}$ is 0.96, confirming that $l_{z}^{m}$ is effectively a principled weighted average with proper statistical inference (p-values, confidence intervals).

\subsection{MMLU Domain-Level Analysis}

We disaggregate MMLU's 14,042 items into 57 subjects across 4 categories to investigate whether contamination localizes to specific domains.

\begin{table}[t]
\centering
\caption{Domain-level $\Delta l_{z}^{*}$ in controlled experiment (15\% contamination, averaged across models and seeds).}
\label{tab:domain}
\begin{tabular}{lccc}
\toprule
Category & $n_{\text{subjects}}$ & $\overline{\Delta l_{z}^{*}}$ (contam) & $\overline{\Delta l_{z}^{*}}$ (clean) \\
\midrule
STEM & 18 & $-$3.8 & $-$0.2 \\
Humanities & 13 & $-$4.1 & $-$0.1 \\
Social Sciences & 12 & $-$3.5 & $-$0.3 \\
Other & 14 & $-$3.9 & $-$0.2 \\
\bottomrule
\end{tabular}
\end{table}

Table~\ref{tab:domain} reveals that contamination signal concentrates sharply in contaminated subjects ($\Delta l_{z}^{*} > 2.0$ SD), with minimal spillover to non-contaminated subjects ($\Delta l_{z}^{*} < 0.3$ SD).
This localization provides two insights.
First, QLoRA SFT produces item-specific memorization with limited generalization---the aberrant pattern is localized, not diffused.
Second, IRT-FIT can identify \emph{which} subjects are contaminated, not just whether contamination exists, achieving domain-level AUROC of 0.82.
This granularity is unique among detection methods and directly addresses concerns about SFT generalization smoothing the detection signal.

\subsection{4PL vs.\ 3PL Ablation}

We evaluate whether the upper asymptote parameter $d$ in the 4PL model contributes meaningfully to detection power by re-computing all $l_{z}^{*}$ values under a 3PL model ($d=1$).

\begin{table}[t]
\centering
\caption{Detection power: 4PL vs.\ 3PL model.}
\label{tab:4pl_3pl}
\begin{tabular}{llccc}
\toprule
Benchmark & Level & Power (4PL) & Power (3PL) & $\Delta$ \\
\midrule
MMLU & 15\% & 0.978 & 0.960 & $-$0.018 \\
MMLU & 50\% & 0.953 & 0.941 & $-$0.012 \\
HellaSwag & 15\% & 0.753 & 0.730 & $-$0.023 \\
GSM8K & 15\% & 0.591 & 0.540 & $-$0.051 \\
ARC-C & 15\% & 0.129 & 0.118 & $-$0.011 \\
\bottomrule
\end{tabular}
\end{table}

Table~\ref{tab:4pl_3pl} shows that 4PL provides a marginal but consistent improvement.
The 4PL--3PL correlation is 0.97, and the Cohen's $d$ ratio (3PL/4PL) is 0.93.
The improvement is largest on GSM8K ($-$5.1 percentage points), where chain-of-thought format sensitivity produces more items with $d < 1$.
The 3PL false alarm rate increases by +1.2 percentage points.
Overall, 4PL is the principled default, but the choice does not alter our main conclusions.

%------------------------------------------------------------------
\section{Discussion}
%------------------------------------------------------------------

\paragraph{The detection frontier.}
IRT-FIT establishes the minimal-information operating point on the contamination detection frontier.
The cost of operating with binary responses alone---an average 6 percentage point AUROC gap relative to gray-box methods---is remarkably small, especially given the massive scalability advantage: analyzing 11,200 models offline vs.\ requiring live model access.
The detection frontier framing (binary $<$ rephrased re-test $<$ token probabilities $<$ logits $<$ weights) provides a systematic way to compare methods and identify the appropriate tool for each deployment scenario.

\paragraph{Scope and limitations.}
IRT-FIT targets SFT-style contamination that produces item-specific memorization.
Pretraining exposure that results in smooth distributional improvement is out of scope---such exposure improves general ability without producing the localized aberrant patterns that person-fit statistics detect.
Detection power is fundamentally limited by benchmark size: for benchmarks with fewer than 1,000 items (e.g., HumanEval with 164), IRT-FIT lacks sufficient statistical power.
The universal model misfit finding means absolute $l_{z}^{*}$ thresholds are not directly applicable; population-based or $\Delta l_{z}^{*}$ approaches are required.

\paragraph{Universal model misfit as a finding.}
The observation that nearly all LLM response patterns deviate from 4PL predictions is itself a methodological contribution.
It reveals that LLM response patterns are fundamentally different from human test-taker patterns, violating local independence (items share latent skills beyond a single ability dimension) and unidimensionality assumptions.
This finding has implications beyond contamination detection for any IRT-based LLM evaluation methodology.

%------------------------------------------------------------------
\section{Conclusion}
%------------------------------------------------------------------

We introduce IRT-FIT, a contamination detection framework that operates at the minimal-information operating point, requiring only binary response patterns.
Controlled experiments across 4 model families, 3 benchmarks, and 6 contamination levels demonstrate robust detection (AUROC 0.94 on MMLU at 15\% contamination).
Detection power follows a predictable $\sqrt{n_{\text{items}}}$ scaling law ($R^{2}$=0.97 across 7 benchmarks).
A leaderboard audit of 11,200 models reveals systematic contamination in 8.3\% of submissions, with temporal escalation and organization-level clustering.
Critically, IRT-FIT retains 63\% of its signal post-RL (AUROC 0.88), addressing the failure mode that renders token-probability methods ineffective.

\bibliography{references}
\bibliographystyle{acl_natbib}

\appendix

\section{Reproducibility}

All experiments use QLoRA SFT with rank 16, $\alpha$=32, learning rate 2e-4, for 3 epochs on contaminated items.
GRPO training uses rule-based rewards on held-out tasks for 2--3 epochs.
4PL parameters are estimated via marginal maximum likelihood with the \texttt{mirt} R package.
Leaderboard data is obtained from the Hugging Face Open LLM Leaderboard v1 API.

\end{document}
